\begin{ex}
Tính: $\frac{1}{2} \times \frac{1}{3} = \dots$
\shortans{1/6}
\loigiai{
$\frac{1 \times 1}{2 \times 3} = \frac{1}{6}$.
}
\end{ex}

\begin{ex}
Tính: $\frac{2}{3} \times \frac{4}{5} = \dots$
\shortans{8/15}
\loigiai{
$\frac{2 \times 4}{3 \times 5} = \frac{8}{15}$.
}
\end{ex}

\begin{ex}
Tính: $\frac{3}{4} \times \frac{1}{2} = \dots$
\shortans{3/8}
\loigiai{
$\frac{3 \times 1}{4 \times 2} = \frac{3}{8}$.
}
\end{ex}

\begin{ex}
Tính: $\frac{2}{5} \times \frac{5}{7} = \dots$ (tối giản).
\shortans{2/7}
\loigiai{
$\frac{2 \times 5}{5 \times 7} = \frac{2}{7}$.
}
\end{ex}

\begin{ex}
Tính: $\frac{3}{8} \times \frac{4}{9} = \dots$ (tối giản).
\shortans{1/6}
\loigiai{
$\frac{12}{72} = \frac{1}{6}$.
}
\end{ex}

\begin{ex}
Tính: $\frac{5}{6} \times \frac{3}{10} = \dots$ (tối giản).
\shortans{1/4}
\loigiai{
$\frac{15}{60} = \frac{1}{4}$.
}
\end{ex}

\begin{ex}
Tính: $\frac{2}{3} \times 6 = \dots$ (kết quả số tự nhiên).
\shortans{4}
\loigiai{
$\frac{2 \times 6}{3} = 4$.
}
\end{ex}

\begin{ex}
Tính: $\frac{3}{4} \times 8 = \dots$ (kết quả số tự nhiên).
\shortans{6}
\loigiai{
$\frac{3 \times 8}{4} = 6$.
}
\end{ex}

\begin{ex}
Tính: $5 \times \frac{2}{5} = \dots$ (kết quả số tự nhiên).
\shortans{2}
\loigiai{
$\frac{5 \times 2}{5} = 2$.
}
\end{ex}

\begin{ex}
Tính: $\frac{4}{7} \times 0 = \dots$
\shortans{0}
\loigiai{
Mọi số nhân với 0 đều bằng 0.
}
\end{ex}

\begin{ex}
Tính: $\frac{5}{9} \times 1 = \dots$
\shortans{5/9}
\loigiai{
$\frac{5}{9}$.
}
\end{ex}

\begin{ex}
Một hình vuông có cạnh dài $\frac{2}{3}$ m. Chu vi hình vuông đó là bao nhiêu mét?
\shortans{8/3}
\loigiai{
$\frac{2}{3} \times 4 = \frac{8}{3}$ m.
}
\end{ex}

\begin{ex}
Diện tích hình vuông ở bài toán trên là bao nhiêu mét vuông?
\shortans{4/9}
\loigiai{
$\frac{2}{3} \times \frac{2}{3} = \frac{4}{9}\text{ m}^2$.
}
\end{ex}

\begin{ex}
Một hình chữ nhật có chiều dài $\frac{4}{5}$ m, chiều rộng $\frac{1}{2}$ m. Diện tích hình chữ nhật đó là bao nhiêu mét vuông? (tối giản).
\shortans{2/5}
\loigiai{
$\frac{4}{5} \times \frac{1}{2} = \frac{4}{10} = \frac{2}{5}\text{ m}^2$.
}
\end{ex}

\begin{ex}
Tính: $\frac{1}{3} \times \frac{3}{4} \times \frac{4}{5} = \dots$ (tối giản).
\shortans{1/5}
\loigiai{
Triệt tiêu được $\frac{1}{5}$.
}
\end{ex}

\begin{ex}
Tính: $\frac{2}{5} \times \frac{3}{4} = \dots$ (tối giản).
\shortans{3/10}
\loigiai{
$\frac{6}{20} = \frac{3}{10}$.
}
\end{ex}

\begin{ex}
Tính nhanh: $\frac{3}{7} \times \frac{5}{8} + \frac{3}{7} \times \frac{3}{8}$.
\shortans{3/7}
\loigiai{
$\frac{3}{7} \times (\frac{5}{8} + \frac{3}{8}) = \frac{3}{7} \times 1 = \frac{3}{7}$.
}
\end{ex}

\begin{ex}
Tính: $12 \times \frac{3}{4} = \dots$
\shortans{9}
\loigiai{
$12 \times 3 : 4 = 9$.
}
\end{ex}

\begin{ex}
Tính: $\frac{5}{8} \times 16 = \dots$
\shortans{10}
\loigiai{
$16 : 8 \times 5 = 10$.
}
\end{ex}

\begin{ex}
Một tấm bìa hình chữ nhật có chiều dài $\frac{3}{4}$ m, chiều rộng $\frac{2}{3}$ m. Diện tích tấm bìa là bao nhiêu mét vuông?
\shortans{1/2}
\loigiai{
$\frac{3}{4} \times \frac{2}{3} = \frac{6}{12} = \frac{1}{2}\text{ m}^2$.
}
\end{ex}

\begin{ex}
Tính: $\frac{7}{9} \times \frac{9}{7} = \dots$
\shortans{1}
\loigiai{
Tích của hai phân số đảo ngược bằng 1.
}
\end{ex}

\begin{ex}
Tính: $\frac{4}{5} \times \frac{5}{4} = \dots$
\shortans{1}
\loigiai{
Bằng 1.
}
\end{ex}

\begin{ex}
Tính: $\frac{1}{2} \times \frac{2}{3} \times \frac{3}{4} = \dots$
\shortans{1/4}
\loigiai{
Triệt tiêu được $\frac{1}{4}$.
}
\end{ex}

\begin{ex}
Tính: $\frac{5}{6} \times 18 = \dots$
\shortans{15}
\loigiai{
$18 : 6 \times 5 = 15$.
}
\end{ex}

\begin{ex}
Tìm $x$: $x : \frac{2}{3} = \frac{3}{4}$.
\shortans{1/2}
\loigiai{
$x = \frac{3}{4} \times \frac{2}{3} = \frac{6}{12} = \frac{1}{2}$.
}
\end{ex}
